\scanpage{80}
aux groupes fondamentaux premiers aux car résiduelles\dots) Par un argument général, le normalisateur de $L_i$ dans $E_U$ s'envoie \uline{sur} $E_S$, on trouve donc des scindages pour cette extension, qu'on peut regarder comme une extension
\[
(13)\quad 1\longrightarrow T\longrightarrow N(L_i)\longrightarrow E_S\longrightarrow1
\]
La classe d'isomorphisme est un élément \deleted{\unsure{de}}
\[
(14)\quad c_i\in H^2(E_S,\underset{\substack{\Vert\\\varprojlim\mu_n}}{T}),
\]
que je me propose d'étudier. On a, si $S$ est un $K(\pi,1)$
\[
(15)\quad H^2(E_S,T)\simeq H^2(S,T)\simeq\varprojlim_{\unsure{n}} H^2(S,\mu_n)
\]
d'ailleurs on a une suite exacte de Kummer (\unsure{où} $\operatorname{Pic}(S)=H^1(S,G_m)$)
\[
(16)\quad 0-\operatorname{Pic}(S)_n-H^2(S,\mu_n)\longrightarrow {}_nH^2(S,G_m)\longrightarrow0
\]
d'où par passage à la limite
\[
(17)\quad 0-
\underset{\substack{\Vert\,\text{déf}\!^{\mathrm n}\\\varprojlim(\operatorname{Pic}_n=\operatorname{Pic}/n\operatorname{Pic})}}{\operatorname{Pic}(S)^\wedge}
\longrightarrow H^2(S,T)\longrightarrow\varprojlim {}_nH^2(S,G_m)
\]
Dans le cas où $S$ est un schéma élémentaire sur un corps de \deleted{\unsure{car.}} \unsure{abst.|de t.}
\margin{79}
